% Curs 12 - Pagina 3 (3-I-1995).
% Fragment pentru % Curs 12 - Pagina 3 (3-I-1995).
% Fragment pentru % Curs 12 - Pagina 3 (3-I-1995).
% Fragment pentru % Curs 12 - Pagina 3 (3-I-1995).
% Fragment pentru \input{courses/pages/curs12-pagina3-ro}.
% Transcriere fidelă în ordinea manuscrisului, fără preambul şi fără salt forţat de pagină.

\subsection*{Curs 12 - Pagina 3}

\noindent\underline{Monoidul liber generat de o mul\c{t}ime}

\bigskip
\noindent Fie $A\neq\varnothing$, $1\notin A$ ($A$ se nume\c{s}te \underline{alfabet},
 elementele lui $A$ se numesc \underline{litere}).

\bigskip
\noindent
\[
A^{+}:=\bigl\{u\mid u:\mathbb{N}^{*}\longrightarrow A\cup\{1\},\ \exists\, n\geq 1
\text{ a.\^{\i}. }
\begin{array}{ll}
 u(i)\in A, & (\forall)\ i=1,\ldots,n,\\
 u(i)=1, & (\forall)\ i>n
\end{array}
\bigr\}.
\]

\noindent $u=(x_1,x_2,\ldots,x_m,1,\ldots)$,
\qquad $v\in A^{+}$, $v=(y_1,\ldots,y_n,1,\ldots)$.

\medskip
\noindent Avem
\[
u=v\quad\Longleftrightarrow\quad m=n\ \text{\c{s}i}\ x_i=y_i
\quad(\forall)\ i=1,\ldots,m=n.
\]

\bigskip
\noindent $A^{+}\ni u,v$,
\[
u\cdot v\ \overset{\mathrm{def}}{=}
(x_1,\ldots,x_m,y_1,\ldots,y_n,1,\ldots).
\]
\noindent\hspace*{3em}\textit{concatenare}\hspace{5em}\textit{juxtapunere}
% Manuscrisul indică aceste două denumiri prin săgeţi sub produs şi şirul rezultat.

\bigskip
\noindent $(A^{+},\cdot)$ este semigrup (asociativ).

\bigskip
\noindent $A^{*}=A^{+}\cup\{1\}$.

\noindent Extindem opera\c{t}ia ``$\cdot$'' la $A^{*}$:
\[
u\cdot 1=1u=1,\qquad (\forall)\ u\in A, \qquad 1\cdot 1=1.
\]
% ATENŢIE: partea finală „=1” a primei relaţii pare a fi în manuscris.
% Pentru element neutru ar trebui „=u”; nu se corectează tacit.
% Şi domeniul „u în A” este redat aşa cum apare, deşi extinderea
% ar cere de regulă u în A^*.

\bigskip
\noindent\underline{Obs.}\quad $(A^{*},\cdot,1)$ este monoid,
numit \underline{monoidul liber generat de $A$}.

\medskip
\noindent $A\subseteq A^{*}$.
% Manuscrisul notează includerea mulţimii de litere în monoid printr-o acoladă.

\noindent
\[
\zeta_A:A\longrightarrow A^{+},\qquad
\zeta_A(x)=(x,1,1,\ldots),\qquad\text{aplica\c{t}ie injectiv\u{a}}.
\]
% Simbolul aplicaţiei pare „zeta_A” în manuscris; de verificat dacă este iota_A.

\medskip
\noindent\underline{Obs.}\quad
$u=(x_1,x_2,\ldots,x_m,1,1,\ldots)
=\zeta_A(x_1)\zeta_A(x_2)\cdots\zeta_A(x_m)$.

\medskip
\noindent Elementele lui $A^{+}$ se numesc \underline{cuvinte peste $A$};
$m$ se nume\c{s}te lungimea cuv\^{a}ntului.

\bigskip
\noindent $(A^{*},\cdot,1)$ este monoid ``tare'' necomutativ:
\[
x,y\in A,\quad x\neq y\quad\Longrightarrow\quad xy\neq yx.
\]

\bigskip
\noindent\fbox{T.}\quad
\underline{Proprietatea de universalitate a semigrupului liber generat de o mul\c{t}ime}

\medskip
\noindent Fie $A\neq\varnothing$,
\[A\xrightarrow{\ \zeta_A\ } A^{+}.\]

\noindent $(\forall)$ semigrup $(S,\cdot)$ \c{s}i
$(\forall)$ aplica\c{t}ie de mul\c{t}imi $\zeta:A\longrightarrow S$,
\[
(\exists!)\ f:A^{+}\longrightarrow S\quad\text{morfism de semigrup}
\]
\noindent a.\^{\i}. $f\circ\zeta_A=\zeta$.

\[
\begin{tikzcd}[column sep=large,row sep=large]
A \arrow[r,"\zeta_A"] \arrow[dr,"\zeta"'] & A^{+} \arrow[d,dashed,"f"] \\
& S
\end{tikzcd}
\]
% Teorema este numai enunţată pe această pagină; nu adăugăm demonstraţie.
.
% Transcriere fidelă în ordinea manuscrisului, fără preambul şi fără salt forţat de pagină.

\subsection*{Curs 12 - Pagina 3}

\noindent\underline{Monoidul liber generat de o mul\c{t}ime}

\bigskip
\noindent Fie $A\neq\varnothing$, $1\notin A$ ($A$ se nume\c{s}te \underline{alfabet},
 elementele lui $A$ se numesc \underline{litere}).

\bigskip
\noindent
\[
A^{+}:=\bigl\{u\mid u:\mathbb{N}^{*}\longrightarrow A\cup\{1\},\ \exists\, n\geq 1
\text{ a.\^{\i}. }
\begin{array}{ll}
 u(i)\in A, & (\forall)\ i=1,\ldots,n,\\
 u(i)=1, & (\forall)\ i>n
\end{array}
\bigr\}.
\]

\noindent $u=(x_1,x_2,\ldots,x_m,1,\ldots)$,
\qquad $v\in A^{+}$, $v=(y_1,\ldots,y_n,1,\ldots)$.

\medskip
\noindent Avem
\[
u=v\quad\Longleftrightarrow\quad m=n\ \text{\c{s}i}\ x_i=y_i
\quad(\forall)\ i=1,\ldots,m=n.
\]

\bigskip
\noindent $A^{+}\ni u,v$,
\[
u\cdot v\ \overset{\mathrm{def}}{=}
(x_1,\ldots,x_m,y_1,\ldots,y_n,1,\ldots).
\]
\noindent\hspace*{3em}\textit{concatenare}\hspace{5em}\textit{juxtapunere}
% Manuscrisul indică aceste două denumiri prin săgeţi sub produs şi şirul rezultat.

\bigskip
\noindent $(A^{+},\cdot)$ este semigrup (asociativ).

\bigskip
\noindent $A^{*}=A^{+}\cup\{1\}$.

\noindent Extindem opera\c{t}ia ``$\cdot$'' la $A^{*}$:
\[
u\cdot 1=1u=1,\qquad (\forall)\ u\in A, \qquad 1\cdot 1=1.
\]
% ATENŢIE: partea finală „=1” a primei relaţii pare a fi în manuscris.
% Pentru element neutru ar trebui „=u”; nu se corectează tacit.
% Şi domeniul „u în A” este redat aşa cum apare, deşi extinderea
% ar cere de regulă u în A^*.

\bigskip
\noindent\underline{Obs.}\quad $(A^{*},\cdot,1)$ este monoid,
numit \underline{monoidul liber generat de $A$}.

\medskip
\noindent $A\subseteq A^{*}$.
% Manuscrisul notează includerea mulţimii de litere în monoid printr-o acoladă.

\noindent
\[
\zeta_A:A\longrightarrow A^{+},\qquad
\zeta_A(x)=(x,1,1,\ldots),\qquad\text{aplica\c{t}ie injectiv\u{a}}.
\]
% Simbolul aplicaţiei pare „zeta_A” în manuscris; de verificat dacă este iota_A.

\medskip
\noindent\underline{Obs.}\quad
$u=(x_1,x_2,\ldots,x_m,1,1,\ldots)
=\zeta_A(x_1)\zeta_A(x_2)\cdots\zeta_A(x_m)$.

\medskip
\noindent Elementele lui $A^{+}$ se numesc \underline{cuvinte peste $A$};
$m$ se nume\c{s}te lungimea cuv\^{a}ntului.

\bigskip
\noindent $(A^{*},\cdot,1)$ este monoid ``tare'' necomutativ:
\[
x,y\in A,\quad x\neq y\quad\Longrightarrow\quad xy\neq yx.
\]

\bigskip
\noindent\fbox{T.}\quad
\underline{Proprietatea de universalitate a semigrupului liber generat de o mul\c{t}ime}

\medskip
\noindent Fie $A\neq\varnothing$,
\[A\xrightarrow{\ \zeta_A\ } A^{+}.\]

\noindent $(\forall)$ semigrup $(S,\cdot)$ \c{s}i
$(\forall)$ aplica\c{t}ie de mul\c{t}imi $\zeta:A\longrightarrow S$,
\[
(\exists!)\ f:A^{+}\longrightarrow S\quad\text{morfism de semigrup}
\]
\noindent a.\^{\i}. $f\circ\zeta_A=\zeta$.

\[
\begin{tikzcd}[column sep=large,row sep=large]
A \arrow[r,"\zeta_A"] \arrow[dr,"\zeta"'] & A^{+} \arrow[d,dashed,"f"] \\
& S
\end{tikzcd}
\]
% Teorema este numai enunţată pe această pagină; nu adăugăm demonstraţie.
.
% Transcriere fidelă în ordinea manuscrisului, fără preambul şi fără salt forţat de pagină.

\subsection*{Curs 12 - Pagina 3}

\noindent\underline{Monoidul liber generat de o mul\c{t}ime}

\bigskip
\noindent Fie $A\neq\varnothing$, $1\notin A$ ($A$ se nume\c{s}te \underline{alfabet},
 elementele lui $A$ se numesc \underline{litere}).

\bigskip
\noindent
\[
A^{+}:=\bigl\{u\mid u:\mathbb{N}^{*}\longrightarrow A\cup\{1\},\ \exists\, n\geq 1
\text{ a.\^{\i}. }
\begin{array}{ll}
 u(i)\in A, & (\forall)\ i=1,\ldots,n,\\
 u(i)=1, & (\forall)\ i>n
\end{array}
\bigr\}.
\]

\noindent $u=(x_1,x_2,\ldots,x_m,1,\ldots)$,
\qquad $v\in A^{+}$, $v=(y_1,\ldots,y_n,1,\ldots)$.

\medskip
\noindent Avem
\[
u=v\quad\Longleftrightarrow\quad m=n\ \text{\c{s}i}\ x_i=y_i
\quad(\forall)\ i=1,\ldots,m=n.
\]

\bigskip
\noindent $A^{+}\ni u,v$,
\[
u\cdot v\ \overset{\mathrm{def}}{=}
(x_1,\ldots,x_m,y_1,\ldots,y_n,1,\ldots).
\]
\noindent\hspace*{3em}\textit{concatenare}\hspace{5em}\textit{juxtapunere}
% Manuscrisul indică aceste două denumiri prin săgeţi sub produs şi şirul rezultat.

\bigskip
\noindent $(A^{+},\cdot)$ este semigrup (asociativ).

\bigskip
\noindent $A^{*}=A^{+}\cup\{1\}$.

\noindent Extindem opera\c{t}ia ``$\cdot$'' la $A^{*}$:
\[
u\cdot 1=1u=1,\qquad (\forall)\ u\in A, \qquad 1\cdot 1=1.
\]
% ATENŢIE: partea finală „=1” a primei relaţii pare a fi în manuscris.
% Pentru element neutru ar trebui „=u”; nu se corectează tacit.
% Şi domeniul „u în A” este redat aşa cum apare, deşi extinderea
% ar cere de regulă u în A^*.

\bigskip
\noindent\underline{Obs.}\quad $(A^{*},\cdot,1)$ este monoid,
numit \underline{monoidul liber generat de $A$}.

\medskip
\noindent $A\subseteq A^{*}$.
% Manuscrisul notează includerea mulţimii de litere în monoid printr-o acoladă.

\noindent
\[
\zeta_A:A\longrightarrow A^{+},\qquad
\zeta_A(x)=(x,1,1,\ldots),\qquad\text{aplica\c{t}ie injectiv\u{a}}.
\]
% Simbolul aplicaţiei pare „zeta_A” în manuscris; de verificat dacă este iota_A.

\medskip
\noindent\underline{Obs.}\quad
$u=(x_1,x_2,\ldots,x_m,1,1,\ldots)
=\zeta_A(x_1)\zeta_A(x_2)\cdots\zeta_A(x_m)$.

\medskip
\noindent Elementele lui $A^{+}$ se numesc \underline{cuvinte peste $A$};
$m$ se nume\c{s}te lungimea cuv\^{a}ntului.

\bigskip
\noindent $(A^{*},\cdot,1)$ este monoid ``tare'' necomutativ:
\[
x,y\in A,\quad x\neq y\quad\Longrightarrow\quad xy\neq yx.
\]

\bigskip
\noindent\fbox{T.}\quad
\underline{Proprietatea de universalitate a semigrupului liber generat de o mul\c{t}ime}

\medskip
\noindent Fie $A\neq\varnothing$,
\[A\xrightarrow{\ \zeta_A\ } A^{+}.\]

\noindent $(\forall)$ semigrup $(S,\cdot)$ \c{s}i
$(\forall)$ aplica\c{t}ie de mul\c{t}imi $\zeta:A\longrightarrow S$,
\[
(\exists!)\ f:A^{+}\longrightarrow S\quad\text{morfism de semigrup}
\]
\noindent a.\^{\i}. $f\circ\zeta_A=\zeta$.

\[
\begin{tikzcd}[column sep=large,row sep=large]
A \arrow[r,"\zeta_A"] \arrow[dr,"\zeta"'] & A^{+} \arrow[d,dashed,"f"] \\
& S
\end{tikzcd}
\]
% Teorema este numai enunţată pe această pagină; nu adăugăm demonstraţie.
.
% Transcriere fidelă în ordinea manuscrisului, fără preambul şi fără salt forţat de pagină.

\subsection*{Curs 12 - Pagina 3}

\noindent\underline{Monoidul liber generat de o mul\c{t}ime}

\bigskip
\noindent Fie $A\neq\varnothing$, $1\notin A$ ($A$ se nume\c{s}te \underline{alfabet},
 elementele lui $A$ se numesc \underline{litere}).

\bigskip
\noindent
\[
A^{+}:=\bigl\{u\mid u:\mathbb{N}^{*}\longrightarrow A\cup\{1\},\ \exists\, n\geq 1
\text{ a.\^{\i}. }
\begin{array}{ll}
 u(i)\in A, & (\forall)\ i=1,\ldots,n,\\
 u(i)=1, & (\forall)\ i>n
\end{array}
\bigr\}.
\]

\noindent $u=(x_1,x_2,\ldots,x_m,1,\ldots)$,
\qquad $v\in A^{+}$, $v=(y_1,\ldots,y_n,1,\ldots)$.

\medskip
\noindent Avem
\[
u=v\quad\Longleftrightarrow\quad m=n\ \text{\c{s}i}\ x_i=y_i
\quad(\forall)\ i=1,\ldots,m=n.
\]

\bigskip
\noindent $A^{+}\ni u,v$,
\[
u\cdot v\ \overset{\mathrm{def}}{=}
(x_1,\ldots,x_m,y_1,\ldots,y_n,1,\ldots).
\]
\noindent\hspace*{3em}\textit{concatenare}\hspace{5em}\textit{juxtapunere}
% Manuscrisul indică aceste două denumiri prin săgeţi sub produs şi şirul rezultat.

\bigskip
\noindent $(A^{+},\cdot)$ este semigrup (asociativ).

\bigskip
\noindent $A^{*}=A^{+}\cup\{1\}$.

\noindent Extindem opera\c{t}ia ``$\cdot$'' la $A^{*}$:
\[
u\cdot 1=1u=1,\qquad (\forall)\ u\in A, \qquad 1\cdot 1=1.
\]
% ATENŢIE: partea finală „=1” a primei relaţii pare a fi în manuscris.
% Pentru element neutru ar trebui „=u”; nu se corectează tacit.
% Şi domeniul „u în A” este redat aşa cum apare, deşi extinderea
% ar cere de regulă u în A^*.

\bigskip
\noindent\underline{Obs.}\quad $(A^{*},\cdot,1)$ este monoid,
numit \underline{monoidul liber generat de $A$}.

\medskip
\noindent $A\subseteq A^{*}$.
% Manuscrisul notează includerea mulţimii de litere în monoid printr-o acoladă.

\noindent
\[
\zeta_A:A\longrightarrow A^{+},\qquad
\zeta_A(x)=(x,1,1,\ldots),\qquad\text{aplica\c{t}ie injectiv\u{a}}.
\]
% Simbolul aplicaţiei pare „zeta_A” în manuscris; de verificat dacă este iota_A.

\medskip
\noindent\underline{Obs.}\quad
$u=(x_1,x_2,\ldots,x_m,1,1,\ldots)
=\zeta_A(x_1)\zeta_A(x_2)\cdots\zeta_A(x_m)$.

\medskip
\noindent Elementele lui $A^{+}$ se numesc \underline{cuvinte peste $A$};
$m$ se nume\c{s}te lungimea cuv\^{a}ntului.

\bigskip
\noindent $(A^{*},\cdot,1)$ este monoid ``tare'' necomutativ:
\[
x,y\in A,\quad x\neq y\quad\Longrightarrow\quad xy\neq yx.
\]

\bigskip
\noindent\fbox{T.}\quad
\underline{Proprietatea de universalitate a semigrupului liber generat de o mul\c{t}ime}

\medskip
\noindent Fie $A\neq\varnothing$,
\[A\xrightarrow{\ \zeta_A\ } A^{+}.\]

\noindent $(\forall)$ semigrup $(S,\cdot)$ \c{s}i
$(\forall)$ aplica\c{t}ie de mul\c{t}imi $\zeta:A\longrightarrow S$,
\[
(\exists!)\ f:A^{+}\longrightarrow S\quad\text{morfism de semigrup}
\]
\noindent a.\^{\i}. $f\circ\zeta_A=\zeta$.

\[
\begin{tikzcd}[column sep=large,row sep=large]
A \arrow[r,"\zeta_A"] \arrow[dr,"\zeta"'] & A^{+} \arrow[d,dashed,"f"] \\
& S
\end{tikzcd}
\]
% Teorema este numai enunţată pe această pagină; nu adăugăm demonstraţie.
